% TOP_PROBLEM: {"id": 30000968, "title": "Generalized Hodge Conjecture", "category_id": 5, "category": {"id": 5, "name": "algebraic_geometry", "display_name": "Algebraic Geometry", "description": "Geometric objects defined by polynomial equations.", "slug": "algebraic-geometry", "order_index": 5, "created_at": "2026-07-31T15:26:25.671Z"}, "status": "open"}
% FIRST_POSTED: null
% !TeX program = lualatex
% Compile twice: lualatex open_problems_500.tex
% All references and prose are embedded; no BibTeX database or external inputs.
\documentclass[11pt,a4paper]{article}
\usepackage[margin=24mm,headheight=14pt]{geometry}
\usepackage{fontspec}
\setmainfont{Latin Modern Roman}
\setsansfont{Latin Modern Sans}
\setmonofont{Latin Modern Mono}
\usepackage{amsmath,amssymb,mathtools}
\usepackage{unicode-math}
\setmathfont{Latin Modern Math}
\usepackage{microtype}
\usepackage{enumitem,xurl,needspace}
\usepackage[dvipsnames]{xcolor}
\usepackage{hyperref}
\hypersetup{unicode=true,colorlinks=true,linkcolor=MidnightBlue,urlcolor=MidnightBlue,
 pdftitle={Mathematical Research Notebook},pdfauthor={Source-annotated compilation},bookmarksnumbered=true}
\usepackage{bookmark}
\usepackage{tocloft}
\setlength{\cftsecnumwidth}{3em}
\cftsetpnumwidth{2.5em}
\cftsetrmarg{4em}
\usepackage{titlesec}
\titleformat{\section}{\Large\bfseries\raggedright}{\thesection}{0.7em}{}
\usepackage{fancyhdr}
\pagestyle{fancy}\fancyhf{}
\fancyhead[L]{\small\sffamily Open Problems Math | Research notebook}
\fancyhead[R]{\small 27 September 2026}
\fancyfoot[C]{\thepage}
\setlength{\parindent}{0pt}
\setlength{\parskip}{4pt plus 1pt}
\setlength{\emergencystretch}{3em}
\setcounter{tocdepth}{1}
\setcounter{secnumdepth}{1}
\setlist[itemize]{leftmargin=3em,itemsep=3pt,topsep=4pt,parsep=0pt}
\allowdisplaybreaks
\newcommand{\N}{\mathbb{N}}
\newcommand{\Z}{\mathbb{Z}}
\newcommand{\Q}{\mathbb{Q}}
\newcommand{\R}{\mathbb{R}}
\newcommand{\C}{\mathbb{C}}
\newcommand{\F}{\mathbb{F}}
\newcommand{\A}{\mathbb{A}}
\newcommand{\PP}{\mathbb P}
\newcommand{\E}{\mathbb{E}}
\newcommand{\eps}{\varepsilon}
\newcommand{\dd}{\,\mathrm d}
\newcommand{\sref}[2]{\hyperlink{source:#1:#2}{[S#2]}}
\newcommand{\eref}[2]{\hyperlink{extra:#1:#2}{[E#2]}}
% Turn the next switch off for a shorter reading copy. The quotations preserve
% every exactTarget field, including qualifications omitted from a short summary.
\newif\ifshowcatalogue
\showcataloguetrue
\newcommand{\cataloguescope}[1]{\ifshowcatalogue
 \paragraph{Exact scope in the supplied catalogue.}
 \begin{quote}\small #1\end{quote}\fi}
\usepackage{etoolbox}
\AtBeginEnvironment{itemize}{\small}
% Standalone entry; original section content is delimited for verification.
\begin{document}
\setcounter{section}{0}
%% BEGIN ORIGINAL SECTION CONTAINER
% ================================================================
% problemId: 30000968
% Canonical problem ID: 30000968
% exactTarget SHA-256: 14fa013aa4ae93b7f6d8711da82b7463b7a662ba56ae4f238a4d325aad98aa43
\clearpage
\section*{\texorpdfstring{Generalized Hodge Conjecture}{Generalized Hodge Conjecture}}\label{30000968}
\subsection{Definitions and mathematical statement}
A rational Hodge substructure $V\subseteq H^k(X,\Q)$ has Hodge coniveau at least $c$ if $V^{a,b}=0$ whenever $a<c$ or $b<c$. Let $M^cH^k$ be the sum of all such substructures. Define geometric coniveau by
\[
 N^cH^k(X,\Q)=\sum_{\substack{Z\subseteq X\text{ closed algebraic}\\ \operatorname{codim}Z\ge c}}
 \ker\bigl(H^k(X,\Q)\to H^k(X\setminus Z,\Q)\bigr).
\]
The conjecture is $M^cH^k(X,\Q)=N^cH^k(X,\Q)$ for every smooth projective complex $X$ and all relevant $k,c$. Finite sums can be supported on the union of their supporting algebraic sets.
\par\smallskip\textit{Formulation and concept references:} \sref{30000968}{1}.
\cataloguescope{For every smooth projective complex variety X and integers k,c, prove that the largest rational Hodge substructure of H\textasciicircum{}k(X,Q) with Hodge coniveau at least c equals the subspace of classes that vanish on X minus some closed algebraic subset of codimension at least c.}
\subsection{Short English statement}
Does every part of cohomology whose Hodge types allow high-codimension support actually have such algebraic support?
%% BEGIN RESEARCH INSERTION
\subsection{Research attempt}
\paragraph{Current frontier.}
\textit{Research check: 27 September 2026.}
The general problem is about rational sub-Hodge structures, not merely
individual Hodge classes. Voisin's 2025 survey is the recorded formulation;
her earlier author-hosted treatment gives the support/Gysin description and
the inclusion $N^c\subseteq M^c$ in Theorem 1.9 and Corollary 1.10.
\sref{30000968}{1}, \eref{30000968}{1}.
The 2026 moduli-space manuscript proves specified degree/genus ranges,
with separate hypotheses for its arithmetic statements; its Theorems A--C
are not a theorem for arbitrary smooth projective varieties. \sref{30000968}{3}.
The calculation below is an independent investigation of the realization
step. No historical novelty is asserted for its constituent observations.

\paragraph{Attack route.}
Three routes suggest themselves: produce a supporting subvariety directly;
realize the effective twist of a Hodge substructure on a smaller-weight
cohomology group and algebraize the resulting correspondence; or obtain
support from a degeneration. The first needs geometric existence, the
second needs both realization and algebraicity, and the third needs a
specialization statement that preserves support codimension. We pursue the
second route because it cleanly separates two often conflated obstructions.

\paragraph{Attempt: a correspondence realization criterion.}
Fix $X$ of dimension $n$, $0\leq c\leq k/2$, and
$V\subseteq H^k(X,\Q)$ of Hodge coniveau at least $c$. Put $m=k-2c$.
Suppose there are a smooth projective $Y$ of dimension $d$, and a morphism
of rational Hodge structures
\[
 u:H^m(Y,\Q)(-c)\longrightarrow H^k(X,\Q),\qquad \operatorname{im}u=V.
 \tag{24.1}
\]
Assume additionally that the specific Hodge class representing $u$ in
$H^{2d+2c}(Y\times X,\Q)$ is algebraic. Then $V\subseteq N^cH^k(X,\Q)$.

Here are all the dimension and support details. Poincare duality identifies
$H^m(Y)^*$ with $H^{2d-m}(Y)(d)$. Consequently $u$, viewed without the
source twist as a map of bidegree $(c,c)$, determines a rational Kunneth
class
\[
 \gamma_u\in H^{2d-m}(Y,\Q)\otimes H^k(X,\Q)
       \subset H^{2d+2c}(Y\times X,\Q)
 \tag{24.2}
\]
of type $(d+c,d+c)$. Its action is
$\alpha\mapsto (p_X)_*(p_Y^*\alpha\smile\gamma_u)$.
The duality identification can be chosen with its usual graded signs, so
that this action is exactly $u$, not just its image up to an unspecified
scalar. Algebraicity supplies a rational cycle
$\Gamma=\sum_i a_i\Gamma_i$ of codimension $d+c$ with class $\gamma_u$.
Every component has dimension $n-c$. Properness of $p_X$ shows that
\[
 Z=\bigcup_i p_X(\Gamma_i)\quad\hbox{is closed},\qquad
 \dim Z\leq n-c.
 \tag{24.3}
\]
On $Y\times(X\setminus Z)$ the cycle is empty. Compatibility of cycle
classes, restriction, and proper pushforward therefore makes the
correspondence action zero after restriction to $X\setminus Z$.
This proves the criterion. It uses one finite supporting union for the
whole image, not a potentially infinite union for individual classes.

A precise sufficient package for the full conjecture is now visible:
every $V(c)$ must be a quotient of some $H^{k-2c}(Y,\Q)$, and the classes
(24.2) must be algebraic. The ordinary Hodge conjecture on all products
would supply the second statement, but the first is a further geometric
realization requirement. Calling the twist ``effective'' supplies only
nonnegative Hodge indices, not such a $Y$.

\paragraph{Attempt: the first two remaining weights.}
For $m=0$, $V(c)$ is a sum of copies of $\Q(0)$, so a finite disjoint union
of points supplies (24.1). Algebraicity of (24.2) is then precisely the
ordinary codimension-$c$ Hodge-class problem on $X$. Thus the criterion
contains, rather than bypasses, the ordinary Hodge conjecture.

For $m=1$, $V(c)$ is polarizable of type $(1,0)+(0,1)$. The classical
weight-one Hodge/abelian-variety correspondence realizes it as $H^1(A,\Q)$
up to isogeny. Equivalently, this realization follows by dualizing a
rational lattice, using the polarization for the Riemann form, and forming
the resulting projective complex torus. A rational complement in
$H^1(A)$, if needed, is obtained from a polarization. \eref{30000968}{2}, Theorem 6.20.
Hence in degree $k=2c+1$ the outstanding step in this route is algebraicity
of a particular codimension $\dim A+c$ class on $A\times X$.
It is not enough that divisors on $A$ are algebraic: that class has a mixed
Kunneth component involving $H^k(X)$.

There is also an unconditional boundary range, worth using as a normalization
check. If $n\leq k\leq2n$, put $r=k-n$ and $m'=2n-k$. A smooth intersection
$i:Y\hookrightarrow X$ of $r$ sufficiently ample divisors has
$\dim Y=m'$. With $h$ their common hyperplane class,
\[
 i_*i^*\alpha=h^r\alpha,\qquad
 h^r:H^{m'}(X,\Q)\xrightarrow{\sim}H^k(X,\Q).
 \tag{24.4}
\]
The second equality is hard Lefschetz. Thus every class in $H^k(X)$ has
support of codimension $r$, and
$N^rH^k=M^rH^k=H^k$. This familiar special case is proved here to test the
support bookkeeping; it says nothing about the larger values $c>r$.

\paragraph{Stress test.}
A rational Hodge projector is not automatically an algebraic projector:
replacing the class (24.2) by a ``cycle representing it'' without the
additional assumption is exactly the missing step. Nor does the
semisimplicity of polarizable Hodge structures manufacture a smooth
projective variety realizing an arbitrary effective twist. The argument
has not deduced that realization from a dimension count.
For $k<2c$, both sides vanish; for $c\leq0$, both are all of $H^k$ under
the usual coniveau conventions. Negative or out-of-range cohomological
degrees are zero. These harmless ranges do not conceal a missing positive
coniveau case. Finally, mod-$2$ or integral coniveau obstructions would not
be counterexamples to this rational statement.

\paragraph{Outcome.}
\textbf{Outcome: Conditional reduction.}
The attempt proves a support-producing criterion with explicit cycle
codimension, isolates realization separately from algebraicity, and checks
the low-weight and high-degree boundary cases. A universal resolution
still requires the two stated inputs. In particular, even the weight-zero
case would resolve the ordinary Hodge conjecture; that dependency is a
warning against treating the correspondence class as already algebraic.
%% END RESEARCH INSERTION
\subsection{Sources}
\begin{itemize}\raggedright
\item[\textnormal{[S1]}]\hypertarget{source:30000968:1}{} Claire Voisin, Hodge and generalized Hodge conjectures, coniveau and algebraic cycles, Journal of Open Mathematical Problems 1(1) (2025), 16-51.
\url{https://jomprob.org/index.php/jomp/article/view/Vol-1Issue-1Paper-2}.
{\small\textit{Catalogue formulation source.}}
\item[\textnormal{[S2]}]\hypertarget{source:30000968:2}{} Incidence equivalence, a survey.
\url{https://arxiv.org/abs/2607.22233}.
{\small\textit{Catalogue research/status reference; not a proof endorsement.}}
\item[\textnormal{[S3]}]\hypertarget{source:30000968:3}{} On the Hodge and Tate conjectures for moduli spaces of curves.
\url{https://arxiv.org/html/2605.20453v1}.
{\small\textit{Catalogue research/status reference; not a proof endorsement.}}
%% BEGIN RESEARCH REFERENCES
\item[\textnormal{[E1]}]\hypertarget{extra:30000968:1}{} Claire Voisin, \emph{Hodge Structures, Coniveau and Algebraic Cycles}, author-hosted manuscript, Theorem 1.9 and Corollary 1.10.
\url{https://webusers.imj-prg.fr/~claire.voisin/Articlesweb/23-.pdf}.
{\small\textit{Added research source. Support, Gysin maps and the coniveau inclusion. The new realization calculation is given in the text. Consulted 27 September 2026.}}
\item[\textnormal{[E2]}]\hypertarget{extra:30000968:2}{} Pierre Deligne and James S. Milne, \emph{Tannakian Categories}, in \emph{Hodge Cycles, Motives, and Shimura Varieties}, Lecture Notes in Mathematics 900 (1982), 101--228; corrected author version dated 15 August 2012, Theorem 6.20.
\url{https://www.jmilne.org/math/xnotes/tc.pdf}.
{\small\textit{Added research source. Weight-one polarizable Hodge structures of abelian-variety type. Consulted 27 September 2026.}}
%% END RESEARCH REFERENCES
\end{itemize}

%% END ORIGINAL SECTION CONTAINER
\end{document}
